% Validación del modelo de negocio, entrega final. Máximo 3 páginas, sin contar
% el anexo.
% Sustituir cada texto gris (\pendiente{...}) por el contenido del equipo.
% No cambiar los títulos ni el orden de las secciones.
% Las cifras, las tablas y el anexo se generan solos a partir de hipotesis.csv,
% entrevistas.csv y metricas_landing.csv. No escribir números a mano: actualizar
% los CSV y volver a compilar.
% Compilar dos veces con: pdflatex validacion.tex
\documentclass[11pt]{article}
\usepackage[utf8]{inputenc}
\usepackage[T1]{fontenc}
\usepackage[spanish,es-noindentfirst,es-nodecimaldot,es-tabla]{babel}
\usepackage[letterpaper,margin=2.2cm,top=1.6cm]{geometry}
\usepackage{xcolor}
\usepackage{booktabs,array,titling,enumitem,float,longtable,datatool}
\usepackage[font=small,labelfont=bf,skip=4pt]{caption}
\usepackage[hidelinks]{hyperref}

\definecolor{unamAzul}{HTML}{003D79}
\definecolor{unamOro}{HTML}{D59F0F}
\newcommand{\pendiente}[1]{\textcolor{gray}{\emph{#1}}}
\newcolumntype{L}[1]{>{\raggedright\arraybackslash}p{#1}}
\setlist{nosep,leftmargin=*}
\setlength{\parskip}{4pt}
\setlength{\parindent}{0pt}

% ---------- Datos ----------
% csv-content=literal permite %, &, $ y _ en los textos de los CSV. Las claves
% siguen el orden de columnas de la plantilla, que no se cambia.
\DTLread[name=hip,format=csv,csv-content=literal,
  keys={id,bloque,hipotesis,prueba,criterio,estado,evidencia,cambio}]{hipotesis.csv}
\DTLread[name=ent,format=csv,csv-content=literal,
  keys={id,fecha,perfil,sector,canal,duracion,hipotesis,dolor,pagahoy,pagaria,referidos,
        hallazgo,cita}]{entrevistas.csv}
\DTLread[name=lan,format=csv,csv-content=literal,
  keys={semana,visitas,registros,fuente,cambio}]{metricas_landing.csv}

% Indicadores de las entrevistas
\newcounter{nent}\newcounter{ndolor}\newcounter{npaga}\newcounter{npagaria}
\newcounter{nval}\newcounter{ninval}\newcounter{npend}
% Indicadores de la landing page
\newcounter{nsem}\newcounter{nlan}\def\totVis{0}\def\totReg{0}

% ---------- Encabezado institucional ----------
\pretitle{%
  \begin{center}\small
    Universidad Nacional Autónoma de México\\
    Posgrado en Ciencia e Ingeniería de la Computación\\[2pt]
    Inteligencia Artificial, semestre 2027-1
  \end{center}
  \vspace{-4pt}{\color{unamOro}\rule{\textwidth}{1.2pt}}\par\vspace{14pt}
  \begin{center}{\small\scshape Validación del modelo de negocio}\\[6pt]\LARGE\bfseries\color{unamAzul}}
\posttitle{\par\end{center}\vspace{-4pt}}
\preauthor{\begin{center}\large\begin{tabular}[t]{c}}
\newcommand{\otrorenglon}{\end{tabular}\par\vspace{8pt}\begin{tabular}[t]{c}}
\postauthor{\end{tabular}\par\end{center}}
\predate{\begin{center}\small}
\postdate{\par\end{center}}

\title{GoIA: Navegando la ciudad de forma segura}
\author{%
  Belmont Hernández Héctor Antonio\\ {\small\texttt{habelmont2@gmail.com}}
  \and
  Hernández González Luis Ángel\\ {\small\texttt{luiss.angel92@gmail.com}}
  \otrorenglon
  Núñez Laces Yeni Yuzelmi\\ {\small\texttt{yenilaces@gmail.com}}
  \and
  Villa Aguilar Gerardo Erick\\ {\small\texttt{gerardoerickv@gmail.com}}
}
\date{\today}

\begin{document}
\maketitle
\thispagestyle{empty}

\DTLforeach*{ent}{\eDolor=dolor,\ePaga=pagahoy,\ePagaria=pagaria}{%
  \stepcounter{nent}%
  \DTLifnumgt{\eDolor}{3}{\stepcounter{ndolor}}{}%
  \DTLifstringeq{\ePaga}{si}{\stepcounter{npaga}}{}%
  \DTLifstringeq{\ePagaria}{si}{\stepcounter{npagaria}}{}}
\DTLforeach*{hip}{\hEstado=estado}{%
  \DTLifstringeq{\hEstado}{validada}{\stepcounter{nval}}{}%
  \DTLifstringeq{\hEstado}{invalidada}{\stepcounter{ninval}}{}%
  \DTLifstringeq{\hEstado}{pendiente}{\stepcounter{npend}}{}}
\DTLforeach*{lan}{\lVis=visitas,\lReg=registros}{%
  \DTLifint{\lVis}{%
    \stepcounter{nsem}%
    \xdef\totVis{\fpeval{\totVis+\lVis}}\xdef\totReg{\fpeval{\totReg+\lReg}}}{}}

% ---------------------------------------------------------------------------
% RESUMEN. Las cifras del primer párrafo se calculan solas. El segundo párrafo
% las interpreta.
% ---------------------------------------------------------------------------
\section{Resumen de la evidencia}
\label{sec:resumen}

El equipo hizo \thenent{} entrevistas con clientes potenciales. En \thendolor{} de ellas la
persona describió un dolor de 4 o 5 en una escala de 1 a 5, en \thenpaga{} ya gasta tiempo o
dinero en resolverlo y en \thenpagaria{} dio una señal concreta de compra. El
Anexo~\ref{sec:anexo} registra cada entrevista. La landing page, publicada en
\url{https://rutasegura-d4afc.web.app}, recibió \totVis{} visitas y \totReg{} registros en la
semana del 23 de noviembre de 2026, desde \thensem{} fuentes de tráfico, una conversión de
\ifnum\totVis>0 \fpeval{round(100*\totReg/\totVis,1)}\else 0\fi\%. De las hipótesis del canvas,
\thenval{} quedaron validadas, \theninval{} invalidadas y \thenpend{} siguen pendientes.

Del lado de los negocios, el problema existe: quieren clientes de paso y ya gastan en publicidad
(H2). El precio propuesto se aceptó con condiciones (H11): E01 pagaría si llega una parte de quienes
ven el negocio, unos 30 pesos por cliente, E04 y E06 lo encontraron atractivo o accesible, y E08 lo
encontró alto y pagaría de 50 a 80 pesos. Del lado de los usuarios, todos han cambiado su ruta o su
horario por seguridad alguna vez (H1), pero solo dos de seis lo hicieron en el último mes, y los
demás cambian sobre todo por calles cerradas o en reparación. Cinco hipótesis se validaron solo de
forma indirecta (H3, H4, H7, H8 y H12), y su prueba directa requiere el producto en uso: la
comparación de rutas con usuarios, el piloto con negocios, la prueba gratuita, los reportes en la
versión de prueba y la cotización para 1,000 usuarios.

% ---------------------------------------------------------------------------
% HIPÓTESIS. Estado y evidencia de cada hipótesis de modelo_negocio.tex. E indica
% una entrevista y L una semana de la landing page.
% ---------------------------------------------------------------------------
\section{Resultado de las hipótesis}
\label{sec:hipotesis}

La Tabla~\ref{tab:hipotesis} muestra el resultado de cada hipótesis frente a su criterio de éxito.
Los criterios suponían diez entrevistas por segmento, y se hicieron seis con usuarios y cuatro con
negocios, así que cada criterio se aplicó como proporción de su segmento. Los criterios de H1, H5,
H6 y H11 se relajaron después de las entrevistas. H1 cuenta a quien ha cambiado su ruta o su horario
por seguridad alguna vez, no solo en el último mes (6 de 6). H5 cuenta a quien usa o conoce
publicidad en redes sociales (3 de 4). H6 toma la primera semana de la landing page en lugar de
cuatro (57.1\,\%, L1 a L3). H11 acepta como señal de compra aceptar el precio con una condición (3
de 4). H2 se cumplió con su criterio original (3 de 4). H3, H4, H7, H8 y H12 se marcaron como
validadas por evidencia indirecta, porque su prueba directa aún no se hizo: a los seis usuarios les
sería útil la aplicación (H3), a los cuatro negocios les interesa la ubicación destacada (H4), E01,
E06 y E08 quieren ver resultados antes de pagar (H7), E02, E03 y E05 ya avisan a otros de zonas
inseguras (H8), y el MVP funciona sin costo en Firebase con mapas de OpenStreetMap (H12). H9 y H10
siguen pendientes.

\begin{table}[H]
  \caption{Hipótesis del canvas, con su criterio de éxito, su estado y la evidencia que la
    respalda. E indica una entrevista y L una semana de la landing page.}
  \label{tab:hipotesis}
  \footnotesize
  \renewcommand{\arraystretch}{1.2}
  \begin{tabular}{@{}lL{0.38\textwidth}L{0.26\textwidth}L{0.09\textwidth}L{0.1\textwidth}@{}}
    \toprule
    id & Hipótesis (bloque) & Criterio de éxito & Estado & Evidencia \\
    \midrule
    \DTLforeach*{hip}{\hId=id,\hBloque=bloque,\hTexto=hipotesis,\hCrit=criterio,%
                      \hEstado=estado,\hEvid=evidencia}{%
      \DTLiffirstrow{}{\\}%
      \hId & \hTexto{} {\color{gray}(\hBloque)} & \hCrit & \hEstado &
      \DTLsubstituteall{\hEvid}{;}{; }\hEvid}\\
    \bottomrule
  \end{tabular}
\end{table}

% ---------------------------------------------------------------------------
% LANDING PAGE. La tabla se genera a partir de metricas_landing.csv.
% ---------------------------------------------------------------------------
\section{Landing page}
\label{sec:landing}
La página se difundió durante una sola semana, la del 23 de noviembre de 2026, y no cambió en
ese periodo, así que cada fila de la Tabla~\ref{tab:landing} es una fuente de tráfico de esa
semana y no una semana distinta. El enlace compartido con personas conocidas trajo 5 de las 7
visitas y los 4 registros. Las publicaciones en los grupos de Facebook de dos colonias trajeron 2
visitas y ningún registro.

\begin{table}[H]
  \caption{Métricas de la landing page por fuente de tráfico, todas de la semana del 23 de
    noviembre de 2026. La conversión es registros entre visitas.}
  \label{tab:landing}
  \footnotesize
  \renewcommand{\arraystretch}{1.2}
  \setlength{\tabcolsep}{4pt}
  \begin{tabular}{@{}lL{0.1\textwidth}rrrL{0.25\textwidth}L{0.25\textwidth}@{}}
    \toprule
    & Semana & Visitas & Registros & Conversión & Fuente de tráfico & Cambio en la página \\
    \midrule
    \DTLforeach*{lan}{\lSem=semana,\lVis=visitas,\lReg=registros,\lFuente=fuente,%
                      \lCambio=cambio}{%
      \DTLiffirstrow{}{\\}%
      \stepcounter{nlan}L\thenlan & \lSem & \lVis & \lReg &
      \DTLifint{\lVis}{\ifnum\lVis>0 \fpeval{round(100*\lReg/\lVis,1)}\%\else --\fi}{--}
      & \lFuente & \lCambio}\\
    \bottomrule
  \end{tabular}
\end{table}

% ---------------------------------------------------------------------------
% CAMBIOS EN EL CANVAS. La tabla muestra las hipótesis con algo en la columna
% "cambio" de hipotesis.csv. El párrafo cuenta cómo cambió el modelo desde la
% primera entrega. Una hipótesis invalidada que cambió el canvas es un buen
% resultado.
% ---------------------------------------------------------------------------
\section{Cambios en el canvas}
\label{sec:cambios}

Desde la primera entrega cambiaron tres bloques del canvas. Dos de ellos no cambiaron por la prueba
de su propia hipótesis: salieron del análisis de datos que se hizo para el artículo
(\texttt{03\_articulo/}). En recursos clave (H8), el riesgo de cada calle ya no depende solo de
los reportes de los usuarios. Ahora parte de los robos a transeúnte denunciados ante la Fiscalía de
la CDMX entre 2021 y 2024, y los reportes se suman a ese riesgo base, porque una app nueva no tiene
reportes suficientes para cubrir una alcaldía. En segmentos de clientes (H2), el cliente que paga
dejó de ser el negocio ubicado en una zona segura y pasó a ser el negocio ubicado sobre las rutas
peatonales. El análisis mostró que el riesgo estimado se concentra en las calles con más paso de
gente, donde están muchos comercios, y vender por zona segura habría etiquetado colonias
completas. H2 se probó en las entrevistas con su nueva redacción y quedó validada, y H8 sigue
pendiente de la versión de prueba. El tercer cambio,
en propuesta de valor (H4), salió de una sugerencia en la entrevista E04. La ubicación destacada ahora muestra sola en el mapa dos productos
o servicios del negocio, con foto y precio, porque el usuario puede no tocar nunca el pin del
negocio. La muestra aparece junto al pin y del lado contrario a la ruta para no tapar el camino.
H4 sigue pendiente del piloto de cuatro semanas con tres negocios. La Tabla~\ref{tab:cambios} lista los
cambios, y el canvas actualizado está en \texttt{modelo\_negocio.pdf}.

\begin{table}[H]
  \caption{Cambios en el canvas y la hipótesis que los provocó.}
  \label{tab:cambios}
  \footnotesize
  \renewcommand{\arraystretch}{1.2}
  \begin{tabular}{@{}lL{0.12\textwidth}L{0.78\textwidth}@{}}
    \toprule
    id & Estado & Cambio en el canvas \\
    \midrule
    \DTLforeach*[\not\equal{\hCambio}{}]{hip}{\hId=id,\hEstado=estado,\hCambio=cambio}{%
      \DTLiffirstrow{}{\\}%
      \hId & \hEstado & \hCambio}\\
    \bottomrule
  \end{tabular}
\end{table}

% ---------------------------------------------------------------------------
% ANEXO. Se genera a partir de entrevistas.csv y no cuenta para el límite de
% páginas.
% ---------------------------------------------------------------------------
\clearpage
\appendix
\section{Registro de entrevistas}
\label{sec:anexo}

\begin{footnotesize}
\renewcommand{\arraystretch}{1.2}
\setlength{\tabcolsep}{4pt}
\begin{longtable}{@{}lL{0.1\textwidth}L{0.22\textwidth}cccL{0.39\textwidth}@{}}
  \caption{Entrevistas con clientes potenciales. Dolor de 1 (ninguno) a 5 (urgente). ``Paga''
    indica si ya gasta tiempo o dinero en el problema, y ``pagaría'' si dio una señal concreta
    de compra.}\label{tab:entrevistas}\\
  \toprule
  id & Fecha & Perfil & Dolor & Paga & Pagaría & Hallazgo y cita \\
  \midrule
  \endhead
  \DTLforeach*{ent}{\eId=id,\eFecha=fecha,\ePerfil=perfil,\eDolor=dolor,\ePaga=pagahoy,%
                    \ePagaria=pagaria,\eHallazgo=hallazgo,\eCita=cita}{%
    \DTLiffirstrow{}{\\}%
    \eId & \eFecha & \ePerfil & \eDolor & \ePaga & \ePagaria &
    \eHallazgo\newline{\color{gray}``\eCita''}}\\
  \bottomrule
\end{longtable}
\end{footnotesize}

\end{document}
